\documentclass[11pt]{book}
\usepackage[paperwidth=145mm,paperheight=200mm,top=18mm,bottom=20mm,inner=20mm,outer=16mm]{geometry}
\usepackage{brumfiel}
% figures10.tex -- Chapter 10 (AREA) figures redrawn in TikZ.
% Proportions read off scans/scan13.pdf p1-p16 (book pp.180-192).
% Every constrained point is constructed, never eyeballed; the stated
% hypotheses are mirrored in tools/constraints/ch10.py.

% ------------------------------------------------------------ Fig 10-1, 10-2
% 10-1: unit-square counting -- 1x1, 2x1, 3x2.
% 10-2: half-unit square inside a unit square; 2.5 x 3 grid.
\newcommand{\FIGXONETWO}{\begin{bkfigure}[\linewidth]
\begin{minipage}{0.52\linewidth}\centering
\begin{tikzpicture}[scale=0.706]
  % (a) single unit square.  Panels (a) and (b) sit on a band 0.5 units above
  % the old one: their "below" labels used to run into the top edge of the
  % 3x2 grid of panel (c), which the two panels overlap in x.
  \draw[fig] (0,3.1) rectangle (1,4.1);
  \node[slab] at (0.5,3.6) {1};
  \node[slab,left] at (0,3.6) {1};  \node[slab,below] at (0.5,3.1) {1};
  % (b) two unit squares side by side
  \draw[fig] (2.4,3.1) rectangle (4.4,4.1);
  \draw[fig] (3.4,3.1) -- (3.4,4.1);
  \node[slab] at (2.9,3.6) {1}; \node[slab] at (3.9,3.6) {1};
  \node[slab,left] at (2.4,3.6) {1}; \node[slab,right] at (4.4,3.6) {1};
  \node[slab,below] at (2.9,3.1) {1}; \node[slab,below] at (3.9,3.1) {1};
  % (c) 3 x 2 grid
  \draw[fig] (0.7,0) rectangle (3.7,2);
  \draw[fig] (1.7,0) -- (1.7,2); \draw[fig] (2.7,0) -- (2.7,2);
  \draw[fig] (0.7,1) -- (3.7,1);
  \foreach \x in {1.2,2.2,3.2}{\node[slab] at (\x,1.5) {1};
                               \node[slab] at (\x,0.5) {1};
                               \node[slab,below] at (\x,0) {1};}
  \node[slab,left] at (0.7,1.5) {1}; \node[slab,left] at (0.7,0.5) {1};
\end{tikzpicture}
\figcap{10--1}
\end{minipage}\hfill
\begin{minipage}{0.46\linewidth}\centering
\begin{tikzpicture}[scale=0.706]
  % (a) half-unit square inside a dashed unit square.  The book quarters the
  % unit square: a dashed mid-vertical and mid-horizontal run right across it,
  % and only the lower-left quarter is inked solid.  The two dashed halves of
  % those medians were missing.
  \draw[aux] (0,1.8) rectangle (2,3.8);
  \draw[aux] (1,2.8) -- (1,3.8);          % upper half of the mid-vertical
  \draw[aux] (1,2.8) -- (2,2.8);          % right half of the mid-horizontal
  \draw[fig] (0,1.8) rectangle (1,2.8);
  \node[slab,left] at (0,3.3) {$\tfrac{1}{2}$};
  \node[slab,left] at (0,2.3) {$\tfrac{1}{2}$};
  % (b) 2.5 wide x 3 tall grid
  \draw[fig] (3.1,0) rectangle (5.6,3);
  \draw[fig] (4.1,0) -- (4.1,3); \draw[fig] (5.1,0) -- (5.1,3);
  \draw[fig] (3.1,1) -- (5.6,1); \draw[fig] (3.1,2) -- (5.6,2);
  \node[slab,left] at (3.1,2.5) {1}; \node[slab,left] at (3.1,1.5) {1};
  \node[slab,left] at (3.1,0.5) {1};
  \node[slab,below] at (3.6,0) {1}; \node[slab,below] at (4.6,0) {1};
  \node[slab,below] at (5.35,0) {$\tfrac{1}{2}$};
\end{tikzpicture}
\figcap{10--2}
\end{minipage}
\end{bkfigure}}

% ---------------------------------------------------------------- Fig 10-3
% 1.4 by 2.3 rectangle, partly chopped into squares 0.1 on a side.
\newcommand{\FIGXTHREE}{\begin{bkfigure}[0.62\linewidth]
\begin{tikzpicture}[scale=2.015]
  \draw[fig] (0,0) rectangle (2.3,1.4);
  \draw[fig] (1,0) -- (1,1.4);            % 1-unit divisions
  \draw[fig] (2,0) -- (2,1.4);
  \draw[fig] (0,1) -- (2.3,1);            % the 0.4 band
  \foreach \y in {1.1,1.2,1.3} \draw[fig] (0,\y) -- (2.3,\y);
  \foreach \x in {2.1,2.2} \draw[fig] (\x,0) -- (\x,1.4);
  \node[slab,left] at (0,1.2) {0.4};  \node[slab,left] at (0,0.5) {1};
  \node[slab,below] at (0.5,0) {1};   \node[slab,below] at (1.5,0) {1};
  \node[slab,below] at (2.15,0) {0.3};
\end{tikzpicture}
\figcap{10--3}
\end{bkfigure}}

% ---------------------------------------------------------------- Fig 10-4
% rectangle sqrt(2) by 1, the part past the whole unit divided into tenths.
% Structure measured off scan13 p3 by tracking every vertical (solid lines
% are 100% inked down the page, dashed ones 76-83%):
%   left edge (solid) -- unit divider at exactly 1 unit (solid) -- five
%   DASHED tenth-marks -- and the right edge, at sqrt(2) units, falls
%   BETWEEN the fourth and fifth marks.  That is the whole point of the
%   figure: it shows 1.4 < sqrt(2) < 1.5, which the text asserts two lines
%   above it.  The old drawing put six dashed marks inside and ended the
%   rectangle on a mark, so the reading was lost.  The top and bottom edges
%   run past the right edge to the last mark, as they do in the book.
\newcommand{\FIGXFOUR}{\begin{bkfigure}[0.62\linewidth]
\begin{tikzpicture}[scale=2.015]
  % one unit = 1.6; the book's tenth-marks are slightly compressed (0.089
  % unit apart) so that the sqrt(2) edge clears them -- keep that.
  \def\uu{1.6}\def\tt{0.1425}\def\ww{2.2627}\def\hh{1.91}% w = sqrt2 * 1.6
  \def\xe{2.3125}%                                        last tenth-mark
  \draw[fig] (0,0) -- (\xe,0);                          % bottom, extended
  \draw[fig] (0,\hh) -- (\xe,\hh);                      % top, extended
  \draw[fig] (0,0) -- (0,\hh);                          % left edge
  \draw[fig] (\ww,0) -- (\ww,\hh);                      % right edge = sqrt2
  \draw[fig] (\uu,0) -- (\uu,\hh);                      % the whole unit
  \foreach \k in {1,2,3,4,5}
    \draw[aux] ({\uu+\k*\tt},0) -- ({\uu+\k*\tt},\hh);
  \node[slab] at (0.81,1.74) {1};
  \node[slab,left] at (0,0.955) {1};
  \node[slab] at (1.42,2.18) {0.1};
  \draw[fig,->] (1.60,2.10) -- (1.70,1.97);             % leader into a strip
  % dimension line under the rectangle; the label sits in a break of it,
  % as the book draws it.
  \draw[fig,|<->|] (0,-0.40) -- node[slab,fill=white,inner sep=1.5pt]
        {$\sqrt{2}$} (\ww,-0.40);
\end{tikzpicture}
\figcap{10--4}
\end{bkfigure}}

% ---------------------------------------------------------------- Fig 10-5
% (a) isosceles triangle ABC and rectangle on the same base;
% (b) parallelogram A'B'C'D' with A'B' = D'E';
% (c) two parallelograms on the base A''D''.
\newcommand{\FIGXFIVE}{\begin{bkfigure}[\linewidth]
\begin{minipage}[b]{0.31\linewidth}\centering
\begin{tikzpicture}[scale=0.62]
  % height read off scan13 p4: h : BC = 570 : 535, i.e. 1.07 (was 0.73).
  \coordinate (B) at (0,0); \coordinate (C) at (3,0);
  \coordinate (E) at ($(B)!0.5!(C)$);
  \coordinate (A) at ($(E)+(0,3.2)$);
  \coordinate (D) at ($(B)+(0,3.2)$);
  \draw[aux] (D) -- (A); \draw[aux] (D) -- (B);
  \draw[fig] (B) -- (C) -- (A) -- cycle;
  \draw[fig] (A) -- (E);
  % the book marks AB and AC with DOUBLE ticks
  \foreach \s/\t in {A/B, A/C}
    {\foreach \f in {0.465,0.535}
       {\draw[tick] ($(\s)!\f!(\t)!2.2pt!90:(\t)$) --
                    ($(\s)!\f!(\t)!2.2pt!-90:(\t)$);}}
  \foreach \p/\n/\d in {A/A/above, B/B/below left, C/C/below right,
                        D/D/above left, E/E/below}
    {\node[vlab,\d] at (\p) {$\n$};}
  \node[slab,below=20pt] at (1.5,0) {(a)};
\end{tikzpicture}
\end{minipage}\hfill
\begin{minipage}[b]{0.31\linewidth}\centering
\begin{tikzpicture}[scale=0.62]
  % re-measured on scan13 p4: A'B' rises at 61 degrees (the old 49 made the
  % parallelogram far too flat), |A'B'| : |A'D'| = 0.54, and A'D' drops 6.7
  % degrees left to right.
  \coordinate (Ap) at (0,0.75); \coordinate (Bp) at (0.675,1.98);
  \coordinate (Dp) at (2.586,0.446);
  \coordinate (Cp) at ($(Dp)+(Bp)-(Ap)$);
  \coordinate (Ep) at ($(Dp)-(Bp)+(Ap)$);
  \draw[fig] (Ap) -- (Bp) -- (Cp) -- (Dp) -- cycle;
  \draw[aux] (Bp) -- (Dp); \draw[aux] (Ap) -- (Ep); \draw[aux] (Ep) -- (Dp);
  \foreach \p/\n/\d in {Ap/A'/left, Bp/B'/above left, Cp/C'/above right,
                        Dp/D'/right, Ep/E'/below}
    {\node[vlab,\d] at (\p) {$\n$};}
  \node[slab,below=20pt] at (1.3,-0.5) {(b)};
\end{tikzpicture}
\end{minipage}\hfill
\begin{minipage}[b]{0.31\linewidth}\centering
\begin{tikzpicture}[scale=0.62]
  % Re-measured on scan13 p4 (vertices, not label boxes): height = 1.46 x
  % base (not 1.94), B'' at 0.42 of the base, E'' at 0.91.  Line weights
  % corrected too: the book inks A''B''C''D'' and the top edge B''..F'',
  % and draws the SECOND parallelogram's slant sides A''E'' and D''F''
  % dashed, plus the cut E''D''.  There is no A''C'' diagonal in the book;
  % the old drawing had one, and had A''E''F''D'' solid.
  \coordinate (Aq) at (0,0);    \coordinate (Dq) at (1.6,0);
  \coordinate (Bq) at (0.67,2.33);
  \coordinate (Cq) at ($(Bq)+(Dq)-(Aq)$);
  \coordinate (Eq) at (1.46,2.33);
  \coordinate (Fq) at ($(Eq)+(Dq)-(Aq)$);
  \draw[fig] (Aq) -- (Bq);            % first parallelogram, left side
  \draw[fig] (Bq) -- (Fq);            % top edge, through E'' and C''
  \draw[fig] (Dq) -- (Cq);            % first parallelogram, right side
  \draw[fig] (Aq) -- (Dq);            % common base
  \draw[aux] (Aq) -- (Eq); \draw[aux] (Dq) -- (Fq);
  \draw[aux] (Eq) -- (Dq);
  \foreach \p/\n/\d in {Aq/A''/below left, Dq/D''/below right,
                        Bq/B''/above left, Cq/C''/above,
                        Eq/E''/above, Fq/F''/above right}
    {\node[vlab,\d] at (\p) {$\n$};}
  \node[slab,below=20pt] at (0.8,0) {(c)};
\end{tikzpicture}
\end{minipage}
\figcap{10--5}
\end{bkfigure}}

% ------------------------------------------------------------ Fig 10-6, 10-7
\newcommand{\FIGXSIXSEVEN}{\begin{bkfigure}[\linewidth]
\begin{minipage}{0.5\linewidth}\centering
\begin{tikzpicture}[scale=0.868]
  % the shared side DE is drawn dead vertical in the book (scan13 p5);
  % the other six vertices are re-measured off that page relative to D and
  % the length of DE.  A had been drawn too close to E, which crowded the
  % two labels; in the book A sits well to the left of E.
  \coordinate (D) at (1.6,0.3);  \coordinate (E) at (1.6,1.75);
  \coordinate (A) at (0.464,2.389); \coordinate (B) at (-0.779,1.378);
  \coordinate (C) at (0.365,-0.197);
  \coordinate (F) at (4.004,2.389); \coordinate (G) at (4.335,-0.347);
  \draw[fig] (A) -- (B) -- (C) -- (D) -- (G) -- (F) -- (E) -- cycle;
  \draw[aux] (D) -- (E);
  % E is lettered straight ABOVE its vertex in the book; anchoring the label
  % "left" put it on the side EA, which runs up and to the left from E.
  \foreach \p/\n/\d in {A/A/above, B/B/left, C/C/below left, D/D/below,
                        E/E/above, F/F/above right, G/G/below right}
    {\node[vlab,\d] at (\p) {$\n$};}
\end{tikzpicture}
\figcap{10--6}
\end{minipage}\hfill
\begin{minipage}{0.48\linewidth}\centering
\begin{tikzpicture}[scale=0.868]
  % a : b = 0.43 on scan13 p5 (was 0.53)
  \coordinate (C) at (0,0); \coordinate (A) at (3,0);
  \coordinate (B) at (0,1.3); \coordinate (Cp) at (3,1.3);
  \draw[aux] (B) -- (Cp) -- (A);
  \draw[fig] (C) -- (A) -- (B) -- cycle;
  \draw[rmark] (0,0.28) -- (0.28,0.28) -- (0.28,0);
  \node[slab,left] at (0,0.65) {$a$};
  \node[slab,below] at (1.5,0) {$b$};
  \foreach \p/\n/\d in {B/B/above left, C/C/below left, A/A/below right,
                        Cp/C'/above right}
    {\node[vlab,\d] at (\p) {$\n$};}
\end{tikzpicture}
\figcap{10--7}
\end{minipage}
\end{bkfigure}}

% ---------------------------------------------------------------- Fig 10-8
% parallelogram ABCD (base b = AD, height h) inside rectangle AECF; EB = x.
\newcommand{\FIGXEIGHT}{\begin{bkfigure}[0.66\linewidth]
% Proportions read off scan13 p6: x : b = 217 : 137 (1.58) and h : b = 143 : 137
% (1.04).  The old drawing had x : b = 0.4 and h : b = 0.57 -- far too little
% slant, which makes the x-strip of the proof almost invisible.
\begin{tikzpicture}[scale=1.105]
  \coordinate (A) at (0,0);   \coordinate (D) at (2,0);
  \coordinate (B) at (3.1,2.05);
  \coordinate (C) at ($(B)+(D)-(A)$);
  \coordinate (E) at (0,2.05); \coordinate (F) at (5.1,0);
  \draw[aux] (E) -- (C); \draw[aux] (A) -- (F); \draw[aux] (E) -- (A);
  \draw[aux] (C) -- (F);
  \draw[fig] (A) -- (B) -- (C) -- (D) -- cycle;
  \draw[rmark] (0,1.77) -- (0.28,1.77) -- (0.28,2.05);
  \draw[rmark] (5.1,0.28) -- (4.82,0.28) -- (4.82,0);
  \node[slab,above] at (1.55,2.05) {$x$};
  \node[slab,below] at (1,0) {$b$};
  \node[slab,right] at (5.1,1.02) {$h$};
  \foreach \p/\n/\d in {E/E/above left, B/B/above, C/C/above right,
                        A/A/below left, D/D/below, F/F/below right}
    {\node[vlab,\d] at (\p) {$\n$};}
\end{tikzpicture}
\figcap{10--8}
\end{bkfigure}}

% ---------------------------------------------------------------- Fig 10-9
% triangle of base b and height h, completed to a parallelogram.
\newcommand{\FIGXNINE}{\begin{bkfigure}[0.58\linewidth]
\begin{tikzpicture}[scale=1.235]
  % re-measured on scan13 p7: apex at 0.56 of the base, height 0.54 x base
  \coordinate (B) at (0,0); \coordinate (C) at (2.6,0);
  \coordinate (A) at (1.46,1.41);
  \coordinate (Bp) at ($(A)+(C)-(B)$);
  \coordinate (H) at ($(B)!(A)!(C)$);
  \draw[aux] (A) -- (Bp) -- (C);
  \draw[fig] (B) -- (C) -- (A) -- cycle;
  \draw[fig] (A) -- (H);
  \node[slab,left] at ($(A)!0.5!(H)$) {$h$};
  \node[slab,below] at (1.75,0) {$b$};
  \foreach \p/\n/\d in {A/A/above left, Bp/B'/above right, B/B/below left}
    {\node[vlab,\d] at (\p) {$\n$};}
\end{tikzpicture}
\figcap{10--9}
\end{bkfigure}}

% ---------------------------------------------------------- Fig 10-10, 10-11
\newcommand{\FIGXTENELEVEN}{\begin{bkfigure}[\linewidth]
\begin{minipage}{0.5\linewidth}\centering
% Ex. Group 10-5 no. 5 states l || m, 3 ft apart, and AB = CD = 5 ft; the book
% draws it to that scale, with C directly over B and the zigzag standing for
% the 1-mile gap from B to C.  The old drawing had AB = 1.4 and CD = 1.3 (not
% congruent) and a 1.07 : 1 gap-to-AB ratio instead of 3 : 5.  One tikz unit
% is now one foot.
\begin{tikzpicture}[scale=0.40]
  \draw[fig] (0,3) -- (11.6,3);  \node[vlab,right] at (11.6,3) {$l$};
  \draw[fig] (0,0) -- (11.6,0);  \node[vlab,right] at (11.6,0) {$m$};
  \coordinate (A) at (0.6,0);  \coordinate (B) at (5.6,0);
  \coordinate (C) at (5.6,3);  \coordinate (D) at (10.6,3);
  \draw[fig] (C) -- (6.05,2.45) -- (5.15,1.95) -- (6.05,1.45)
             -- (5.15,0.95) -- (5.9,0.45) -- (B);
  \foreach \p in {A,B,C,D} {\dt{\p}}
  \node[vlab,above] at (C) {$C$};  \node[vlab,above] at (D) {$D$};
  \node[vlab,below] at (A) {$A$};  \node[vlab,below] at (B) {$B$};
\end{tikzpicture}
\figcap{10--10}
\end{minipage}\hfill
\begin{minipage}{0.47\linewidth}\centering
\begin{tikzpicture}[scale=0.48]
  \coordinate (A) at (0,0);  \coordinate (D) at (4,0);
  \coordinate (B) at (4,3);  \coordinate (C) at ($(B)+(D)-(A)$);
  \draw[fig] (A) -- (B) -- (C) -- (D) -- cycle;
  \draw[fig] (B) -- (D);
  \draw[rmark] (4,0.4) -- (3.6,0.4) -- (3.6,0);
  \draw[rmark] (4,2.6) -- (4.4,2.6) -- (4.4,3);
  % "above left" grazed AB; step 0.7 off the midpoint along AB's normal
  \node[slab] at ($(A)!0.5!(B)!0.55!90:(B)$) {5};
  \node[slab,below] at ($(A)!0.5!(D)$) {4};
  \foreach \p/\n/\d in {A/A/below left, D/D/below, B/B/above left,
                        C/C/above right}
    {\node[vlab,\d] at (\p) {$\n$};}
\end{tikzpicture}
\figcap{10--11}
\end{minipage}
\end{bkfigure}}

% ---------------------------------------------------------- Fig 10-12, 10-13
\newcommand{\FIGXTWELVETHIRTEEN}{\begin{bkfigure}[\linewidth]
\begin{minipage}{0.5\linewidth}\centering
\begin{tikzpicture}[scale=0.734]
  \coordinate (B) at (0,0);    \coordinate (C) at (4.5,0);
  \coordinate (A) at (1,3);    \coordinate (D) at (3.5,3);
  \draw[fig] (B) -- (C) -- (D) -- (A) -- cycle;
  \coordinate (T) at (2.25,3);
  \draw[fig] (T) -- (2.25,0);
  \draw[aux] (A) -- (C);
  % plain "above" left the label's box touching the top side; lift it clear
  \node[slab,above] at ($($(A)!0.5!(D)$)+(0,0.28)$) {$b' = 10$};
  \node[slab,below] at ($(B)!0.5!(C)$) {$b = 18$};
  % the book prints "h = 12" INSIDE the trapezoid, in the panel between the
  % left side and the altitude, not out in the margin (scan13 p8)
  \node[slab] at (1.40,1.5) {$h = 12$};
  \foreach \p/\n/\d in {A/A/above left, D/D/above right, B/B/below left,
                        C/C/below right}
    {\node[vlab,\d] at (\p) {$\n$};}
\end{tikzpicture}
\figcap{10--12}
\end{minipage}\hfill
\begin{minipage}{0.47\linewidth}\centering
\begin{tikzpicture}[scale=0.568]
  % measured on scan13 p8: b' : b = 0.563, h : b = 0.72 (the old drawing had
  % 0.5 and 0.5).  Both b' and b - b' are printed INSIDE the trapezoid just
  % above the base, and the b of the dimension line sits in a break of it.
  \coordinate (B) at (0,0);   \coordinate (C) at (6,0);
  \coordinate (A) at (1.2,4.3);
  \coordinate (D) at ($(A)+(3.38,0)$);
  \coordinate (P) at (3.38,0);
  \draw[fig] (B) -- (C) -- (D) -- (A) -- cycle;
  \draw[aux] (D) -- (P);
  \node[slab,above] at ($($(A)!0.5!(D)$)+(0,0.30)$) {$b'$};
  \node[slab] at ($($(B)!0.5!(P)$)+(0,0.50)$) {$b'$};
  \node[slab] at ($($(P)!0.5!(C)$)+(0.15,0.50)$) {$b - b'$};
  \draw[fig,|<->|] (0,-1.3) -- node[slab,fill=white,inner sep=1.5pt]
        {$b$} (6,-1.3);
\end{tikzpicture}
\figcap{10--13}
\end{minipage}
\end{bkfigure}}

% --------------------------------------------------------------- Fig 10-14
\newcommand{\FIGXFOURTEEN}{\begin{bkfigure}[0.52\linewidth]
\begin{tikzpicture}[scale=1.118]
  % base : height = 1.75 and lean : height = 0.26 on scan13 p9
  \coordinate (B) at (0,0); \coordinate (C) at (3,0);
  \coordinate (A) at (0.45,1.72);
  \coordinate (D) at ($(A)+(C)-(B)$);
  \draw[fig] (A) -- (B) -- (C) -- (D) -- cycle;
  \draw[fig] (A) -- (C); \draw[fig] (B) -- (D);
  \coordinate (O) at (intersection of A--C and B--D);
  % "above right" put the label along diagonal BD.  The book sets O straight
  % above the crossing, in the wedge between OA and OD.
  \node[vlab,above] at ($(O)+(0,0.16)$) {$O$};
  \foreach \p/\n/\d in {A/A/above left, D/D/above right, B/B/below left,
                        C/C/below right}
    {\node[vlab,\d] at (\p) {$\n$};}
\end{tikzpicture}
\figcap{10--14}
\end{bkfigure}}

% --------------------------------------------------------------- Fig 10-15
\newcommand{\FIGXFIFTEEN}{\begin{bkfigure}[0.52\linewidth]
\begin{tikzpicture}[scale=1.170]
  % apex at 0.30 of the base, h : b = 0.50 (scan13 p9)
  \coordinate (A) at (0,0); \coordinate (C) at (3.4,0);
  \coordinate (B) at (1.0,1.7);
  \draw[fig] (A) -- (C) -- (B) -- cycle;
  \node[slab,above left]  at ($(A)!0.5!(B)$) {$c$};
  \node[slab,above right] at ($(B)!0.5!(C)$) {$a$};
  \node[slab,below]       at ($(A)!0.5!(C)$) {$b$};
  \foreach \p/\n/\d in {A/A/below left, C/C/below right, B/B/above}
    {\node[vlab,\d] at (\p) {$\n$};}
\end{tikzpicture}
\figcap{10--15}
\end{bkfigure}}

% ---------------------------------------------------------- Fig 10-16, 10-17
\newcommand{\FIGXSIXTEENSEVENTEEN}{\begin{bkfigure}[\linewidth]
\begin{minipage}{0.5\linewidth}\centering
\begin{tikzpicture}[scale=0.873]
  % Re-measured off scan13 p9 at 150 dpi (vertices A(232,65) D(700,82)
  % B(145,342) C(620,335)): h : BC = 0.58 and lean : h = 0.31.  An earlier
  % pass had 0.78 / 0.20, which stretched the parallelogram upright.
  \coordinate (B) at (0,0);   \coordinate (C) at (3.5,0);
  \coordinate (A) at (0.64,2.04);
  \coordinate (D) at ($(A)+(C)-(B)$);
  \coordinate (E) at ($(B)!0.25!(D)$);     % BE : ED = 2 : 6
  \draw[fig] (A) -- (B) -- (C) -- (D) -- cycle;
  \draw[fig] (B) -- (D);
  \draw[aux] (E) -- (C);
  \foreach \p/\n/\d in {A/A/above left, D/D/above right, B/B/below left,
                        C/C/below right, E/E/above left}
    {\node[vlab,\d] at (\p) {$\n$};}
\end{tikzpicture}
\figcap{10--16}
\end{minipage}\hfill
\begin{minipage}{0.48\linewidth}\centering
\begin{tikzpicture}[scale=0.786]
  \coordinate (A) at (0,0); \coordinate (B) at (4.333,0);
  \coordinate (C) at (0.641,1.538);        % AC = 5, CB = 12, AB = 13 (/3)
  \coordinate (D) at ($(A)!(C)!(B)$);
  \draw[fig] (A) -- (B) -- (C) -- cycle;
  \draw[fig] (C) -- (D);
  % No right-angle box: the book draws neither the one at C nor the one at D
  % here, though it marks right angles freely in 10-7, 10-8, 10-11 and 10-20.
  % Both facts are stated in the exercise text instead (scan13 p9).
  \foreach \p/\n/\d in {A/A/below left, B/B/right, C/C/above, D/D/below}
    {\node[vlab,\d] at (\p) {$\n$};}
\end{tikzpicture}
\figcap{10--17}
\end{minipage}
\end{bkfigure}}

% ---------------------------------------------------------- Fig 10-18, 10-19
% 10-18: squares on the three sides of a right triangle.
% 10-19: the (a+b) square dissected into four right triangles and a c-square.
\newcommand{\FIGXEIGHTEENNINETEEN}{\begin{bkfigure}[\linewidth]
\begin{minipage}{0.5\linewidth}\centering
\begin{tikzpicture}[scale=1.187]
  \coordinate (R) at (0,0); \coordinate (S) at (0,1.0); \coordinate (T) at (1.4,0);
  \draw[fig] (R) -- (S) -- (T) -- cycle;
  \draw[fig] (R) -- (-1.0,0) -- (-1.0,1.0) -- (S);          % square 1
  \draw[fig] (R) -- (0,-1.4) -- (1.4,-1.4) -- (T);          % square 2
  \coordinate (N) at (1.0,1.4);                              % normal (s1,s2)
  \draw[fig] (S) -- ($(S)+(N)$) -- ($(T)+(N)$) -- (T);       % square 3
  \node[slab] at (-0.5,0.5) {1};
  \node[slab] at (0.7,-0.7) {2};
  \node[slab] at (1.2,1.2)  {3};
\end{tikzpicture}
\figcap{10--18}
\end{minipage}\hfill
\begin{minipage}{0.48\linewidth}\centering
\begin{tikzpicture}[scale=0.736]
  % In the book (scan13 p11, close-up) the SHORT piece of every side is the
  % one lettered a and the long piece is b: going clockwise from the top-left
  % corner the top edge reads a then b, the right edge a then b, the bottom
  % edge b then a, the left edge b then a.  The letters were in the right
  % places here but the two lengths were interchanged, which mirrored the
  % inner square's tilt.  Swapping the values fixes both at once.
  \def\aa{1.6}\def\bb{2.4}\def\ss{4.0}
  \coordinate (P1) at (\bb,0);   \coordinate (P2) at (\ss,\bb);
  \coordinate (P3) at (\aa,\ss); \coordinate (P4) at (0,\aa);
  \draw[fig] (0,0) rectangle (\ss,\ss);
  \draw[fig] (P1) -- (P2) -- (P3) -- (P4) -- cycle;
  \node[slab,below] at (\bb/2,0)        {$b$};
  \node[slab,below] at ({(\bb+\ss)/2},0){$a$};
  \node[slab,right] at (\ss,\bb/2)      {$b$};
  \node[slab,right] at (\ss,{(\bb+\ss)/2}) {$a$};
  \node[slab,above] at (\aa/2,\ss)      {$a$};
  \node[slab,above] at ({(\aa+\ss)/2},\ss) {$b$};
  \node[slab,left]  at (0,\aa/2)        {$a$};
  \node[slab,left]  at (0,{(\aa+\ss)/2}){$b$};
  \node[slab] at (2.83,1.45) {$c$};
  \node[slab] at (2.55,2.83) {$c$};
  \node[slab] at (1.17,2.55) {$c$};
  \node[slab] at (1.45,1.17) {$c$};
\end{tikzpicture}
\figcap{10--19}
\end{minipage}
\end{bkfigure}}

% --------------------------------------------------------------- Fig 10-20
% Bhaskara's dissection: square BAED of side c, four right triangles
% (legs a, b) round a central square of side a - b.  DF || AC, EG || BC.
\newcommand{\FIGXTWENTY}{\begin{bkfigure}[0.56\linewidth]
\begin{tikzpicture}[scale=1.495]
  \coordinate (B) at (0,0);   \coordinate (A) at (3,0);
  \coordinate (E) at (3,3);   \coordinate (D) at (0,3);
  \coordinate (C)  at (1.92,1.44);     % right angle at C, BC = a, CA = b
  \coordinate (Cp) at (1.56,1.92);
  \coordinate (G)  at (1.08,1.56);
  \coordinate (F)  at (1.44,1.08);
  \draw[aux] (B) -- (D) -- (E) -- (A);
  \draw[fig] (B) -- (A);
  \draw[fig] (B) -- (C) -- (A);
  % The central square is NOT inked as a solid quadrilateral in the book
  % (scan13 p14, close-up): three of its sides fall along lines already drawn
  % -- GF lies on DF, C'G lies on EG, and FC lies on the solid leg BC -- and
  % only CC' is added, dashed.  Drawing the square solid put ink on the page
  % that the source does not have and broke the given/auxiliary convention.
  \draw[aux] (D) -- (F); \draw[aux] (E) -- (G);
  \draw[aux] (C) -- (Cp);
  % right-angle marks at G and at C', as the book has them
  \draw[rmark] (1.256,1.692) -- (1.388,1.516) -- (1.212,1.384);
  \draw[rmark] (1.384,1.788) -- (1.516,1.612) -- (1.692,1.744);
  % "a" stepped 0.25 units off BC along its normal (it was grazing the leg
  % at 1.2pt); the book likewise floats it clear, above the midpoint of BC.
  \node[slab] at (0.81,0.92) {$a$};
  \node[slab] at (2.62,0.95) {$b$};
  \node[slab,below] at (1.5,0) {$c$};
  % C is lettered up and to the RIGHT of the vertex in the book; a plain
  % "right" anchor slid the letter down onto the leg CA (0.5pt clearance).
  \foreach \p/\n/\d in {B/B/below left, A/A/below right, D/D/above left,
                        E/E/above right, C/C/above right, F/F/below right}
    {\node[vlab,\d] at (\p) {$\n$};}
  \node[vlab] at (1.38,2.10) {$G$};
\end{tikzpicture}
\figcap{10--20}
\end{bkfigure}}

% --------------------------------------------------------------- Fig 10-21
% quadrilateral ABCD; DE drawn parallel to the diagonal AC, E on line BC.
\newcommand{\FIGXTWENTYONE}{\begin{bkfigure}[0.6\linewidth]
\begin{tikzpicture}[scale=1.235]
  % A and D re-measured on scan13 p14 against BC: A at (0.38, 0.68) and
  % D at (0.97, 0.83) in units of BC.  The quadrilateral had been drawn
  % noticeably taller than the book's.
  \coordinate (B) at (0,0);   \coordinate (C) at (2.2,0);
  \coordinate (A) at (0.84,1.50); \coordinate (D) at (2.13,1.83);
  \coordinate (Dp) at ($(D)+(C)-(A)$);
  \coordinate (Xf) at (6,0);
  \coordinate (E) at (intersection of D--Dp and B--Xf);
  % the base is drawn as the RAY BC extended, arrowhead at E (scan13 p14)
  \draw[fig,->] (B) -- (E);
  \draw[fig] (B) -- (A) -- (D) -- (C);
  \draw[aux] (A) -- (C); \draw[aux] (A) -- (E); \draw[aux] (D) -- (E);
  \foreach \p/\n/\d in {A/A/above left, D/D/above, B/B/below left,
                        C/C/below, E/E/below right}
    {\node[vlab,\d] at (\p) {$\n$};}
\end{tikzpicture}
\figcap{10--21}
\end{bkfigure}}

% ---------------------------------------------------------- Fig 10-22, 10-23
\newcommand{\FIGXTWENTYTWOTHREE}{\begin{bkfigure}[\linewidth]
\begin{minipage}{0.5\linewidth}\centering
\begin{tikzpicture}[scale=1.048]
  % vertices measured off scan13 p15 and rescaled on AB: the book's D sits
  % almost straight above A and C just outside B, giving a much broader top
  % than the drawing had.
  \coordinate (A) at (0.9,0);   \coordinate (B) at (2.6,0);
  \coordinate (C) at (2.74,2.52); \coordinate (D) at (0.96,2.78);
  \coordinate (E) at (0.08,1.29);
  \draw[fig] (A) -- (B) -- (C) -- (D) -- (E) -- cycle;
  \foreach \p/\n/\d in {A/A/below left, B/B/below right, C/C/above right,
                        D/D/above left, E/E/left}
    {\node[vlab,\d] at (\p) {$\n$};}
\end{tikzpicture}
\figcap{10--22}
\end{minipage}\hfill
\begin{minipage}{0.48\linewidth}\centering
\begin{tikzpicture}[scale=0.839]
  \coordinate (B) at (0,0); \coordinate (C) at (4,0);
  \coordinate (A) at (2,3.464);            % equilateral
  \coordinate (P) at (1.9,1.0);
  \coordinate (L) at ($(B)!(P)!(C)$);
  \coordinate (M) at ($(A)!(P)!(C)$);
  \coordinate (N) at ($(A)!(P)!(B)$);
  \draw[fig] (A) -- (B) -- (C) -- cycle;
  \draw[aux] (P) -- (A); \draw[aux] (P) -- (B); \draw[aux] (P) -- (C);
  \draw[aux] (P) -- (L); \draw[aux] (P) -- (M); \draw[aux] (P) -- (N);
  \foreach \p/\n/\d in {A/A/above, B/B/below left, C/C/below right,
                        L/L/below, M/M/right, N/N/left}
    {\node[vlab,\d] at (\p) {$\n$};}
  % the book sets P in the wedge between PL and PC; at our trim the label has
  % to sit further out along that wedge's bisector to clear both dashed lines
  \node[vlab] at (2.19,0.54) {$P$};
\end{tikzpicture}
\figcap{10--23}
\end{minipage}
\end{bkfigure}}

% --------------------------------------------------------------- Fig 10-24
\newcommand{\FIGXTWENTYFOUR}{\begin{bkfigure}[0.46\linewidth]
\begin{tikzpicture}[scale=1.066]
  % scan13 p16, deskewed by the 4.1 degree page rotation: apex at 0.56 of
  % the base, height 0.72 x base (it had been a symmetric 0.5 / 0.5).
  \coordinate (A) at (2.13,2.74); \coordinate (B) at (0,0);
  \coordinate (C) at (3.8,0);
  \draw[fig] (A) -- (B) -- (C) -- cycle;
  \draw[aux] (A) -- ($(B)!0.5!(C)$);
  \draw[aux] (B) -- ($(A)!0.5!(C)$);
  \draw[aux] (C) -- ($(A)!0.5!(B)$);
\end{tikzpicture}
\figcap{10--24}
\end{bkfigure}}


% ---- local macros (hoist into book preamble) ----
% (No \Xstarnote here: ch10 has starred exercises but the book prints no star
% footnote on any of its pages; ch02:311 owns the book-wide first use.)
% Unnumbered "limit" footnote on book p.182 — same reasoning as ch02's
% \IIstarnote: \starnote in brumfiel.sty emits a bare \footnotetext, which
% prints the footnote counter, and the book marks this note with "*" and no
% number.
\newcommand{\Xlimitnote}{\begingroup\renewcommand{\thefootnote}{}%
  \footnotetext{*The word ``limit'' has a precise mathematical sense, which
  need not concern us now. Roughly, the idea is that the areas of the
  approximating rectangles get close to a number that we call the \emph{area}
  of the given rectangle.}\endgroup}
% Figure-name glyphs. The book writes "area (<glyph>ABCD)". In the running
% proofs it uses a plain square for both rectangles and parallelograms; the
% Review of Chapter 10 summary distinguishes four glyphs, so the parallelogram
% and trapezoid are drawn here rather than faked with \square.
\newcommand{\Xrect}{\ensuremath{\square}}
\newcommand{\Xpgram}{\mathord{\tikz[baseline=-0.05ex,line width=0.5pt]{%
  \draw (0,0) -- (1.35ex,0) -- (1.75ex,1.05ex) -- (0.40ex,1.05ex) -- cycle;}}}
\newcommand{\Xtrap}{\mathord{\tikz[baseline=-0.05ex,line width=0.5pt]{%
  \draw (0,0) -- (1.85ex,0) -- (1.45ex,1.05ex) -- (0.42ex,1.05ex) -- cycle;}}}
% Narrow multicol columns plus this chapter's long primed/subscripted math
% atoms ($A''B''C''D''$, $\sg{A_1B_1}/\sg{A_2B_2}$) leave TeX no legal break
% point. Give it a little slack rather than re-wording the book.
\setlength{\emergencystretch}{2.5em}
% ---- end local macros ----

\begin{document}
\setcounter{chapter}{10}

\chapter*{\raggedright\Huge AREA}
\addcontentsline{toc}{chapter}{10.\ Area}
\markboth{AREA}{}
\vspace{-1em}\noindent\rule{\linewidth}{0.6pt}\vspace{1em}
\figurekey

%========================================================= 10-1
\section*{10--1\quad INTRODUCTION}
\addcontentsline{toc}{section}{10--1\quad Introduction}

You already have an intuitive understanding of the concept of \emph{area}. In
this chapter we will describe \emph{area} in much the same way as we described
length in Chapter 5. There we determined a number that we called the length of
a line segment by choosing an arbitrary \emph{unit segment} and
\emph{counting}. Here, for a given geometric figure, we determine a number
that we call the \emph{area} of the figure by first choosing an arbitrary
\emph{unit square} and then \emph{counting}.

It is much more difficult to describe the \emph{counting process} for area
than it is for length. We content ourselves with a crude description of the
counting process as it relates to finding areas of rectangles. Once the area
formula for rectangles is determined, we employ it to derive area formulas for
other simple polygons.

%========================================================= 10-2
\section*{10--2\quad AREA BY COUNTING}
\addcontentsline{toc}{section}{10--2\quad Area by counting}

We understand by \emph{polygon}, of course, the set of points on the sides.
When we refer to the \emph{area} of a \emph{polygon}, we really are considering
the set of interior points. It would be more descriptive to speak of the
\emph{area} of this \emph{interior set}. But for brevity, we simply speak of
the area of a polygon. First, we study rectangular polygons.

We need to explain what we mean by the area of a rectangle. It will help to
consider a few examples.

Having chosen a unit length, a square 1 unit of length on a side is said to
have 1 unit of area. Clearly, if we had two squares side by side, as in
Fig.~10--1, we would want to count the area as 2. Suppose we had a rectangle
with sides of lengths 3 and 2. Figure 10--1 shows that we would want the area
to be 6 units.

If we had a square of side $\tfrac{1}{2}$ unit, then placing it in a unit
square as in Fig.~9--2,\footnote{\emph{Transcriber's note.} The 1961 printing
reads ``Fig.~9--2'' here; the figure meant is plainly Fig.~10--2. Reproduced
as printed.} we see that the area should be only $\tfrac{1}{4}$ unit.

If we had a rectangle with sides $2\tfrac{1}{2}$ and 3 units, what would the
area be?

\FIGXONETWO

\exercisegroup{10--1}
% No star footnote here. Text review 2026-08-17 checked the foot of book p.181
% at 200 dpi in BOTH sources (photo PDF 195 and scan13 p2): the page ends with
% Fig. 10-3 and carries no footnote. ch02:311 owns the book-wide first use of
% the "*Starred exercises ... are optional" note. Do not re-add it here.
\begin{exercisecols}
\begin{exlist}
\item Draw a rectangle whose sides are 3 and 5 units. Determine its area by
  decomposing it into unit squares and counting.
\item Sketch a figure to illustrate why the area of an $m$ by $n$ rectangle is
  $mn$ units. Here, $m$ and $n$ are positive integers.
\item A rectangle is $3\tfrac{1}{3}$ units long and 3 units wide. Find its
  area. What is the largest size square, such that the rectangle can be
  covered by an integral number of squares? Suppose that the rectangle were
  $3\tfrac{1}{3}$ by 2. What is the answer then?
\item[\stex 4.] Suppose that the length of a rectangle is $a$, and its width
  is $b/c$, where $a$, $b$, and $c$ are positive integers. Obtain the area by
  dividing the rectangle into squares $1/c$ units on a side and then counting.
\item[5.] A rectangle has sides 2.2 and 3.3 in. Obtain its area by dividing it
  into small squares and counting.
\item[6.] A rectangle has sides 3.24 and 4.83 ft. Describe how to obtain the
  area by counting squares.
\item[\stex 7.] The sides of a rectangle are $a/b$ and $c/b$ units, where $a$,
  $b$, and $c$ are positive integers. Show how its area is obtained by counting
  squares.
\end{exlist}
\end{exercisecols}

%========================================================= 10-3
\section*{10--3\quad DEFINITION OF RECTANGLE AREA}
\addcontentsline{toc}{section}{10--3\quad Definition of rectangle area}

In the previous section we saw how to obtain the area of a rectangle by
counting squares if the rectangle has sides whose lengths are integers.
Remember that, so far, we have not explicitly defined what area is, but have
used only our previous experience.

\FIGXTHREE

Suppose the lengths of the sides of a rectangle are $l$ and $w$, where $l$ and
$w$ are terminating decimals. If $l$ and $w$ end in 10ths, then we chop up the
rectangle into squares 0.1 on a side. If $l$ and $w$ end in 100ths, we chop the
rectangle into squares 0.01 on a side, and so on. In Fig.~10--3, we have
started to chop up the rectangle, which is 1.4 by 2.3. Is it necessary to draw
more of the small squares in order to count them?

It is rather obvious that the method of counting could be applied to a
rectangle whose sides are any rational numbers.

\exercise{10--2}
\noindent Find, by counting, the area of the rectangle whose sides are
(a)~2.5 and 3.3, (b)~2.23 and 1.77, (c)~5/3 and 7/3, (d)~5/3 and 1/2,
(e)~5/3 and 3/4.

\medskip
In all our problems so far, we have found, by counting squares, that the area
of the rectangle is $lw$.

Now suppose that at least one of the sides of the rectangle is not a rational
number, that is, this side and the unit are not commensurable. To be definite,
let us consider a rectangle whose length is $\sqrt{2}$ and whose width is 1.
(Remember that $\sqrt{2}$ is not rational, and that \emph{approximately}
$\sqrt{2} = 1.4142$.) Examination of Fig.~10--4 shows that the area of the
rectangle is greater than 1.4 and less than 1.5. We may see this by counting
squares. If we were to subdivide the length into 100ths, we would also see
that the area is greater than 1.41 and less than 1.42. As we approximate to
$\sqrt{2}$ more and more closely by rational numbers, for example, by 1.4,
1.41, 1.414, 1.4142,\dots, we will get rectangles whose areas are ever closer
to $\sqrt{2}$.

% Fig. 10-4 sits at the foot of book p.182, i.e. AFTER Exercise 10-2 and after
% the incommensurable-case prose that references it. It was ahead of the
% exercise here; restored to source order.
\FIGXFOUR

Therefore, the situation in which we find ourselves is as follows. If the
sides of the rectangle are rational numbers, we can get the area by counting.
If the sides are not rational, then simple counting does not give us the area.
In this latter case, we must approximate to the sides by rational numbers and
consider the \emph{limit}*\Xlimitnote{} of the areas of the rectangles whose
sides have rational numbers for lengths. It is possible to prove that this
limit exists and is equal to the product $lw$.

You see that to define the area of a rectangle by starting with counting may be
quite troublesome. Therefore, in this book we do not give a detailed definition
of area of a rectangle, but only state the following theorem, which tells us
how to calculate the area.

\begin{theorem}
The area of a rectangle whose sides have lengths $l$ and $w$ is $lw$.
\end{theorem}

\exercisegroup{10--3}
\begin{exercisecols}
\begin{exlist}
\item A square has side $x$ yards. What is its area in square feet?
\item What is the area of a rectangle $3\sqrt{2}$ units long and $2\sqrt{2}$
  units wide?
\end{exlist}
\end{exercisecols}

%========================================================= 10-4
\section*{10--4\quad AREA OF POLYGONS}
\addcontentsline{toc}{section}{10--4\quad Area of polygons}

So far, the only polygons whose areas we have studied are rectangles. In
discussing areas of other polygons, and in particular triangles, perhaps we
should define area in terms of some kind of counting process. If we did, we
would be forced to concern ourselves with limits, so we will determine areas of
polygons by a method used by the Greek geometers. Our basic formula for the
area of a rectangle will enable us to find areas of triangles, and other
figures, without any counting of squares.

Remember that although the Greeks used numbers in commerce, they had no numbers
other than integers in their geometry. They did not compute numerical measures
of angles, lines, or areas. Just as in a primitive society a hunter might
compare his son's height with that of the son of his friend by marking his
spear, so the Greeks devised a method for comparing the sizes of polygons. Of
course, the Greeks had no word in their geometry corresponding to our term
``area,'' since area is a \emph{number}. Instead of saying that two geometric
figures had \emph{equal areas}, they merely remarked that the figures were
\emph{equal}. The Greeks showed that two polygons were \emph{equal} by cutting
them up into congruent triangles. Note that the word \emph{equal} meant
\emph{the same size} to the Greeks. This is different from the word's meaning
in modern mathematics, where \emph{equal} means \emph{identical}. Figure
10--5(a) illustrates what is meant. The Greek geometers would call rectangle
$ADBE$ and isosceles triangle $ABC$ \emph{equal} because
$\tri{ABE} \cong \tri{ABE}$ and $\tri{ABD} \cong \tri{ACE}$.

\FIGXFIVE

\exercise{10--4}
\noindent In Fig.~10--5(b), $A'B'C'D'$ is a parallelogram, and
$\sg{A'B'} \cong \sg{D'E'}$. Show that $A'B'D'E'$ and $A'B'C'D'$ are
\emph{equal} in the sense of Greek geometry. Do the same for parallelograms
$A''B''C''D''$ and $A''E''F''D''$ in Fig.~10--5(c).

\medskip
We use the methods of the Greek geometers to develop area formulas. But before
doing so, we need a few theorems about areas of geometric figures. These
theorems can be proved, using properties of numbers and our geometric
postulates, but their proofs are long and difficult. Accordingly, we treat them
as postulates even though proofs can be given.

\begin{theorem}
Congruent triangles have equal areas.
\end{theorem}

\begin{theorem}
If two simple polygons have one side and only the points of this side in
common, and if, moreover, they have no interior points in common, then the
simple polygon formed by the sides that remain after the common side is
deleted has an area equal to the sum of the areas of the original polygons.
\end{theorem}

Theorem 10--2 is intuitively obvious, since any method of counting squares to
get the area of one triangle can be duplicated for the other triangle. Draw a
figure to illustrate this. Figure 10--6 will make the meaning of Theorem 10--3
clear. The theorem tells us that the area of polygon $ABCDGFE$ is the sum of
the areas of polygons $ABCDE$ and $DEFG$.

\FIGXSIXSEVEN

With these two theorems and our rectangle formula, we can establish theorems on
the areas of triangles and parallelograms.

\begin{theorem}
If $a$ and $b$ are the lengths of the legs of a right triangle, then the
area of the triangle is $\frac{1}{2}ab$.
\end{theorem}

\noindent\emph{Proof.} Let the right triangle be $\tri{ABC}$, as in
Fig.~10--7, and let $C'$ be chosen as indicated in the figure. (How is $C'$
chosen?)

\begin{steps}
\step{$\tri{ABC} \cong \tri{BAC'}$.}{Why?}
\step{Area $(\tri{ABC})$ = area $(\tri{ABC'})$.}{Theorem 10--2.}
\step{Area $(\tri{ABC})$ $+$\newline area $(\tri{ABC'})$ $=$\newline
      area $(\Xrect AC'BC) = ab$.}{Theorem 10--3 and formula for area of
      rectangle.}
\step{2 area $(\tri{ABC}) = ab$.}{Statements 2 and 3.}
\step{Area $(\tri{ABC}) = \frac{1}{2}ab$.}{Statement 4.}
\end{steps}

In a parallelogram, if we call the length of one side the \emph{base}, then the
corresponding \emph{altitude} (or \emph{height}) is the perpendicular distance
from this side to the line containing the opposite side. From Theorem 10--4, we
can deduce the area of a parallelogram.

\begin{theorem}
The area of a parallelogram is equal to the product of its base and its
height.
\end{theorem}

\FIGXEIGHT

\noindent\emph{Proof.} Let the parallelogram be $ABCD$, as in Fig.~10--8, its
base $b = \sg{AD}$, and its height $h$. Consider, as in the figure, the
rectangle $AECF$. (How can the vertices $E$ and $F$ be found?) Let
$\sg{EB} = x$.

\begin{steps}
\step{$\tri{AEB} \cong \tri{CFD}$.}{Why?}
\step{Area $(\tri{AEB})$ = area $(\tri{CFD})$ = $\frac{1}{2}hx$.}{Why?}
\step{Area $(\Xrect AECF)$ $=$\newline area $(\tri{AEB})$ $+$\newline
      area $(\Xrect ABCD)$ $+$\newline area $(\tri{CFD})$.}{Why?}
\step{$\sg{AF} = \sg{EC} = b + x$.}{Why?}
\step{$h(b + x) = \frac{1}{2}hx$ $+$\newline area $(\Xrect ABCD)$ $+$
      $\frac{1}{2}hx$.}{Statements 3 and 4.}
\step{$hb + hx$ = area $(\Xrect ABCD)$ $+$ $hx$.}{Statement 5 and algebra.}
\step{Area $(\Xrect ABCD) = bh$.}{Statement 6 and algebra.}
\end{steps}

We can now quite easily derive the familiar formula,
$A = \frac{1}{2}b \cdot h$, for computing the area of a triangle. Here, $b$
(base) is the length of a side of a triangle, and $h$ (height) is the
perpendicular distance from the opposite vertex to the line containing the
base.

\FIGXNINE

\begin{theorem}
The area of a triangle is given by the formula
$A = \frac{1}{2}b \cdot h$ (Fig.~10--9).
\end{theorem}

The proof is left to the student.

From these last two theorems, we may conclude that parallelograms (or
triangles) with equal bases and heights have equal areas. The area formula for
triangles provides us with a tool for determining areas of polygons in general,
for the area of any polygon can be computed as the sum of the areas of certain
triangles.

% Exercise Group 10-5 carries eight figures (10-10 ... 10-17).  Each one now
% sits inline, immediately after the exercise that names it, and the group
% runs full measure rather than in two columns -- a figure set in a 52 mm
% column would print smaller than the book's own.  (Caleb, 2026-08-17: "for
% the exercise groups, the figures were not placed inline with the text like
% they should be.")  The list was in two \begin{exlist} blocks with an
% explicit \setcounter across the old figure dump; it is now one list, so the
% numbering runs 1-38 by itself.
\exercisegroup{10--5}
\begin{exercisecols}
\begin{exlist}
\item A triangle has an area of 30 and a base of 6. Determine its height.
\item Two adjacent sides of a parallelogram have lengths of 12 and 8 in. These
  sides determine a $30^\circ$ angle. Compute the area of the parallelogram.
\item A triangle has area 50 and altitude 8. What is the length of its base?
\item A rhombus has one angle of $45^\circ$ and is 2 ft on a side. Find its
  area.
\item In Fig.~10--10, $l$ and $m$ are parallel and 3 ft apart.
  $\sg{AB} = \sg{CD} = 5$ ft. If the distance from $B$ to $C$ is 1 mi, what is
  the area of the quadrilateral $ABDC$?
\exfig{\FIGXTEN}
\item What is the area of $\Xrect ABCD$ if the sides are as shown in
  Fig.~10--11 and $BD \perp AD$?
\exfig{\FIGXELEVEN}
\item The area of a triangle is $A$, and its base is $x$. Find the altitude in
  terms of $A$ and $x$.
\item The base of a parallelogram can be expressed in terms of its area and
  height. What is this formula?
\item Two triangles have the same side for base. The height of one triangle
  exceeds that of the other by 4 in. The areas differ by 40 in$^2$. What is the
  length of the common base?
\item The midpoints of the sides of a parallelogram determine a polygon.
  Compare the area of this polygon to the area of the original parallelogram.
\item[\stex 11.] The midpoints $A_1$, $B_1$, and $C_1$ of the sides of
  $\tri{ABC}$ are joined to form a second triangle. Consider also the midpoints
  $A_2$, $B_2$, and $C_2$ of the sides of triangle $A_1B_1C_1$. This process is
  continued. Compare the area of $\tri{A_5B_5C_5}$ with that of $\tri{ABC}$.
\item[\stex 12.] Triangles $ABC$ and $A'B'C'$ are similar. The lengths of two
  corresponding sides are $s$ and $s'$, where $s = 5s'$. If $A$ and $A'$
  represent the respective areas of the triangles, prove that $A = 25A'$.
  Generalize this problem, and establish the theorem that \emph{if two
  triangles are similar, the ratio of their areas is equal to the ratio of the
  squares of the lengths of two corresponding sides.}
\item[13.] $\tri{ABC}$ is similar to $\tri{A'B'C'}$, and $AB = 2A'B'$. How do
  their areas compare? Draw a figure.
\item[\stex 14.] Plans for a house are drawn, so that $\frac{1}{4}$ in. on the
  plan corresponds to 1 ft in the house. If the plans cover an area of
  105 in$^2$, how many square feet in the house?
\item[15.] Two equilateral triangles have sides of 8 and 10 in.,
  respectively. How do their areas compare?
\item[16.] A rhombus has sides 10 in.\ long. If one angle is $30^\circ$, what
  is its area? if $60^\circ$? if $120^\circ$?
\item[17.] Find the area of an isosceles right triangle if one leg is 20 ft
  long.
\item[18.] In Fig.~10--12, $AD \pll BC$. Compute the area of the figure.
\exfig{\FIGXTWELVE}
\item[\stex 19.] Generalize Exercise 18, designating the lengths of the
  parallel sides, or bases, by $b$ and $b'$ and the height by $h$. Of course,
  you are developing a formula for the area of a trapezoid.
\item[20.] Besides the proof suggested by the dotted line in Fig.~10--12,
  Fig.~10--13 also suggests a proof. Here, it is assumed that $b > b'$. Make a
  proof.
\exfig{\FIGXTHIRTEEN}
% Items 11-20 carry explicit labels, which do not advance exlisti, so the
% automatic counter is still at 10 here.  The old two-block layout restored it
% with a \setcounter on the second \begin{exlist}; merging the blocks into one
% list means the same correction has to be made in place.
\setcounter{exlisti}{20}
\item The bases of a trapezoid are 3 and 12, and the altitude is 6. What is
  its area?
\item The area of a trapezoid of height 6 in.\ is 39 in$^2$. If one base is
  6 in.\ what is the other?
\item One base of a trapezoid is twice as long as the other. If the area is
  30 ft$^2$ and the altitude is 5 ft, what are the two bases?
\item An equilateral triangle has sides 18 in.\ long. Find its area in square
  feet.
\item An equilateral triangle has an area of $(9\sqrt{3})/4$. How long is its
  side?
\item In Fig.~10--14, $ABCD$ is a parallelogram. Prove that the four small
  triangles, $\tri{AOB}$, $\tri{BOC}$, etc., all have equal areas.
\exfig{\FIGXFOURTEEN}
\item The areas of two similar triangles are 360 and 250 square units. If a
  side of the first triangle is 8 units, how long is the corresponding side of
  the second triangle?
\item[\stex 28.] Prove that the ratio of the areas of similar polygons is equal
  to the square of the ratio of corresponding sides.
\item[29.] Find the area of a regular hexagon 6 in.\ on a side. What will be
  the area of a regular hexagon whose side is 8 in.?
\item[\stex 30.] If $a$, $b$, and $c$ are the lengths of the sides of any
  triangle (Fig.~10--15), and $s = \frac{1}{2}(a + b + c)$, then the area is
  given by
  \[ A = \sqrt{s(s\!-\!a)(s\!-\!b)(s\!-\!c)}. \]
  This is ``Heron's formula,'' named after an Alexandrian Greek by that name of
  about 100 B.C., possibly later. Show that this formula is correct.
\exfig{\FIGXFIFTEEN}
\item[31.] The sides of a triangle are 8, 11, and 13 in. Find its area.
\item[32.] The sides of a triangle are 17, 25, and 26 ft. Find its area. Find
  the height to the 25-ft side.
\item[33.] Find the area of an isosceles triangle if the base is 12 in.\ and
  the sides 10 in.
\item[34.] Find the area of an equilateral triangle if one side is 6 in.
\item[35.] If, in Fig.~10--16, $ABCD$ is a parallelogram, with $\sg{BE} = 2$
  and $\sg{ED} = 6$, what is the number
  \[ \frac{\text{Area } (BEC)}{\text{Area } (ABCD)}\,? \]
\exfig{\FIGXSIXTEEN}
\item[36.] A square has a diagonal of length 12 in. Compute the area of the
  square. Determine a formula for the area of a square in terms of the length
  $d$ of its diagonal.
\item[37.] The diagonals of a rhombus are 10 in.\ and 6 in., respectively. What
  is its area? Generalize this problem by formulating and proving a theorem.
\item[38.] In Fig.~10--17, $\ang{C}$ is right, $CD \perp AB$, $\sg{AC} = 5$,
  and $\sg{CB} = 12$. Compute the areas of triangles $ABC$, $ACD$, and $BCD$.
  Check your work by verifying that the sum of the area of the latter two
  triangles is equal to the area of $\tri{ABC}$.
\exfig{\FIGXSEVENTEEN}
\end{exlist}
\end{exercisecols}

%========================================================= 10-5
\section*{10--5\quad AREA PROOF OF THE\\PYTHAGOREAN THEOREM}
\addcontentsline{toc}{section}{10--5\quad Area proof of the Pythagorean theorem}

The Greek geometers thought of the Pythagorean theorem, not as a numerical
relationship between lengths of the sides of a right triangle, but rather as a
relationship between areas. Viewed in this light, the theorem may be stated as:

\begin{theorem}
For every right triangle, it is true that the square on the hypotenuse has
an area equal to the sum of the areas of the squares on the other two sides
(see Fig.~10--18).
\end{theorem}

\FIGXEIGHTEENNINETEEN

\begin{center}
Area 1 $+$ area 2 $=$ area 3.
\end{center}

Probably more different proofs have been given for the Pythagorean theorem than
for any other theorem in mathematics. Our earlier proof did not refer to area
at all. It was essentially an ``algebraic style'' proof. We give a proof below
that is a kind of algebraic and area proof. If the sides of the right triangle
are $a$, $b$, and $c$, consider Fig.~10--19.

\noindent\emph{Proof.} Evidently, the area of the large square, of side
$a + b$, is equal to the area of the square of side $c$ plus the areas of the
four triangles. That is,
\[ (a + b)^2 = c^2 + 4(\tfrac{1}{2}ab). \]
But $(a + b)^2 = a^2 + 2ab + b^2$, so that we get
\[ a^2 + 2ab + b^2 = c^2 + 2ab. \]
From algebra then,
\[ a^2 + b^2 = c^2. \]

There are many more proofs of this theorem. Your teacher can show you some if
you are interested.

% multicols dropped: this group carries Fig 10-20, which now sits inline after
% the exercise (no. 8) that names it.
\exercisegroup{10--6}
\begin{exercisecols}
\begin{exlist}
\item Triangles $ABC$ and $A'B'C'$ are similar. $\sg{AB} = 3$,
  $\sg{A'B'} = 5$, and area $(ABC) = 18$. Compute area $(A'B'C')$.
\item Determine the area of an equilateral triangle if the length of one side
  is 10.
\item Show that the area of an equilateral triangle is given by the formula
  $A = (s^2/4) \cdot \sqrt{3}$, where $s$ is the length of a side.
\item Develop a formula for the area of a regular hexagon if each side has
  length $c$.
\item The lengths of two adjacent sides of a parallelogram are 12 in.\ and
  8 in. The included angle is $45^\circ$. Determine the area of the
  parallelogram.
\item The altitude of a rectangle is 12 ft. A diagonal has length 20 ft.
  Compute the area of the rectangle.
\item[\stex 7.] A square and a regular hexagon have equal perimeters. Find the
  ratio of their areas.
\item[\stex 8.] It is said that the Hindu mathematician Bhaskara, born about
  1100 A.D., discovered a proof of the Pythagorean theorem and presented the
  figure like that in Fig.~10--20 with the single word, ``Behold.'' Give a
  proof. In Fig.~10--20 $DF \pll AC$, and $EG \pll BC$.
  [\emph{Hint:} $\sg{FC} = a - b$.]
\exfig{\FIGXTWENTY}
\end{exlist}
\end{exercisecols}

%========================================================= 10-6
\section*{*10--6\quad CONSTRUCTION PROBLEMS}
\addcontentsline{toc}{section}{*10--6\quad Construction problems}

There are many RC constructions related to area concepts. It is interesting to
note that it is possible to construct a triangle whose area is equal to the
area of any given polygon. We indicate in Fig.~10--21 a procedure for
constructing a triangle equal in area to a quadrilateral. $DE$ has been drawn
parallel to diagonal $AC$. Show that area $(ABE)$ = area $(ABCD)$.
[\emph{Hint:} Show that area $(ADC)$ = area $(AEC)$.]

\FIGXTWENTYONE

% multicols dropped: this group carries Figs 10-22 and 10-23, now inline after
% exercises 1 and 10 respectively.
\exercisegroup{10--7}
\begin{exercisecols}
\begin{exlist}
\item Construct a triangle whose area is equal to that of pentagon $ABCDE$
  (Fig.~10--22). Use the method employed above to construct a triangle equal in
  area to a given quadrilateral.
\exfig{\FIGXTWENTYTWO}
\item Show how to construct a square whose area is equal to the sum of the
  areas of two given squares.
\item Solve Exercise 2 with ``sum'' replaced by ``difference.''
\item Construct a square whose area is $\frac{3}{4}$ of the area of a given
  square.
\item Construct a square whose area is equal to the area of a given rectangle.
\item Construct a rectangle whose area is equal to the area of a given
  triangle.
\item Construct a square whose area is the same as the area of a given
  triangle.
\item Construct a triangle similar to a given triangle and having exactly twice
  the area of the given triangle.
\item Show how to construct a polygon of $n - 1$ sides whose area is the same
  as that of an $n$-gon. Assume $n > 3$.
\item[\stex 10.] Prove that if $P$ is any point in the interior of an
  equilateral triangle (Fig.~10--23), the sum of the lengths of the
  perpendiculars from $P$ to the sides of the triangle is equal to the length
  of an altitude of the triangle. [\emph{Hint:} Consider the areas of triangles
  $APB$, $APC$, and $BPC$.]
\exfig{\FIGXTWENTYTHREE}
\end{exlist}
\end{exercisecols}

%========================================================= review
\section*{REVIEW OF CHAPTER 10}
\addcontentsline{toc}{section}{Review of Chapter 10}

\begin{exlist}
\item In this chapter we have studied methods for finding the areas of
  rectangles, triangles, and polygons. All area formulas were derived from the
  formula for the area of a rectangle.
  \begin{flushleft}
  Area $\Xrect$ $=$ length $\cdot$ width $=$ $lw$.\\
  Area $\Xpgram$ $=$ base $\cdot$ altitude $=$ $bh$.\\
  Area $\triangle$ $=$ $\frac{1}{2}$ base $\cdot$ altitude $=$
    $\frac{1}{2}bh$.\\
  Area $\Xtrap$ $=$ $\frac{1}{2}$ altitude $\cdot$ sum of bases $=$
    $h(b + b')/2$.
  \end{flushleft}
\item Review how the above area formulas were obtained.
\item Areas of similar figures are proportional to the squares of corresponding
  sides.
\end{exlist}

\par\medskip\begin{center}\bfseries Review Exercises\end{center}\par\nopagebreak
\begin{exlist}
\item $\tri{A_1B_1C_1} \sim \tri{A_2B_2C_2}$. If area $(A_2B_2C_2) = 72$, and
  $\sg{A_1B_1}/\sg{A_2B_2} = \tfrac{3}{4}$, find area $(A_1B_1C_1)$.
\item Prove that a median of a triangle divides the triangle into two equal
  areas.
\item[\stex 3.] Prove that the medians divide a triangle into six triangles of
  equal areas (Fig.~10--24).
\exfig{\FIGXTWENTYFOUR}
\item[4.] Prove that the area of a square is half the area of a square whose
  side is the diagonal of the given square.
\item[5.] Two similar triangles have areas of 108 and 48 ft$^2$. The first has
  altitude $7\frac{1}{2}$ ft; find the altitude of the second.
\item[6.] Prove that the area of a trapezoid is equal to the altitude times the
  length of the segment joining the midpoints of the nonparallel sides.
\item[7.] A triangle has one angle of $30^\circ$. Prove that the area is
  $\frac{1}{4}$ of the product of the lengths of the sides adjacent to the
  $30^\circ$ angle.
\item[8.] In $\tri{ABC}$, if $D$ is the foot of the perpendicular from $C$ to
  $AB$, prove that $\sg{AC}^2 - \sg{BC}^2 = \sg{AD}^2 - \sg{BD}^2$.
\item[9.] The hypotenuse of a right triangle is 41 ft. One leg is 40 ft. Find
  the length of the other leg.
\item[10.] The median of a trapezoid is 8 in., and one of the bases is 12 in.
  How long is the other base? If the area is $\frac{1}{2}$ft$^2$, what is the
  altitude?
\item[11.] In a right triangle, the altitude from the right angle is 14 in.
  From the foot of the altitude to a vertex is 8 in. How long is the
  hypotenuse?
\end{exlist}

\par\medskip\begin{center}\bfseries Review Test\end{center}\par\nopagebreak
\begin{exlist}
\item What is the area of a parallelogram with sides 10 and 8 and an included
  angle of $30^\circ$?
\item What is the area of a rhombus with side 10 if an angle is $45^\circ$?
\item Find the area of an isosceles right triangle whose legs are 20.
\item One of the bases of a trapezoid is 15, and its height is 6. Find the
  other base if the area is 75.
\item Two similar triangles, $T$ and $T'$, have two corresponding sides 6 and
  15, respectively. What is the ratio of the area of $T$ to the area of $T'$?
\item The length of a diagonal of a square is 15. Find the area.
\item Two similar polygons have areas of 250 and 160. If the length of a side
  of the first is 6, what is the length of the corresponding side of the
  second?
\item Two triangles are similar, and the lengths of the sides of the first
  triangle are 3, 5, and 6. If the length of the greatest side of the other
  triangle is 15, compute the lengths of the other sides.
\item If the legs of a right triangle are in the ratio of 5 to 4, find the
  ratio of the segments formed on the hypotenuse by the altitude to the
  hypotenuse.
\item If one leg of a right triangle has length 25, and the altitude to the
  hypotenuse is 7, what is the area of the triangle?
\end{exlist}

\end{document}
